\documentclass[11pt]{article}
\usepackage[T1]{fontenc}
\usepackage{lmodern,amsmath,amssymb,amsthm,mathtools,microtype,aliascnt}
\usepackage[margin=1in]{geometry}
\usepackage[colorlinks=true,linkcolor=blue,citecolor=blue,urlcolor=blue]{hyperref}
\usepackage[nameinlink,noabbrev]{cleveref}
\newtheorem{theorem}{Theorem}
\newaliascnt{lemma}{theorem}
\newtheorem{lemma}[lemma]{Lemma}
\aliascntresetthe{lemma}
\crefname{lemma}{lemma}{lemmas}
\Crefname{lemma}{Lemma}{Lemmas}
\newaliascnt{proposition}{theorem}
\newtheorem{proposition}[proposition]{Proposition}
\aliascntresetthe{proposition}
\crefname{proposition}{proposition}{propositions}
\Crefname{proposition}{Proposition}{Propositions}
\theoremstyle{remark}
\newaliascnt{remark}{theorem}
\newtheorem{remark}[remark]{Remark}
\aliascntresetthe{remark}
\crefname{remark}{remark}{remarks}
\Crefname{remark}{Remark}{Remarks}
\hypersetup{pdftitle={Weak existence for adaptive porous flow with a negative drag jump},pdfauthor={Kaiwen Jin},pdfsubject={Weak solutions with free transition drag},pdfkeywords={porous flow, weak existence, negative drag jump}}
\newcommand{\R}{\mathbb R}
\newcommand{\Tr}{\operatorname{Tr}}
\newcommand{\Div}{\operatorname{div}}
\newcommand{\supp}{\operatorname{supp}}
\newcommand{\ind}{\mathbf 1}
\newcommand{\norm}[2]{\lVert #1\rVert_{#2}}
\newcommand{\pair}[2]{\int_\Omega \big(#1\big)\cdot #2}
\title{Weak existence for adaptive porous flow with a negative drag jump}
\author{Kaiwen Jin\\\small\href{mailto:yui.ui.siki@gmail.com}{yui.ui.siki@gmail.com}}
\date{October 13, 2026}
\begin{document}
\maketitle

\begin{abstract}
We prove weak existence for a stationary porous-flow model whose drag law
changes at speed one and has a negative jump. At the threshold the drag is
an arbitrary vector field, as in the free-transition formulation of the model.
The result holds in every dimension $d\ge2$ on bounded connected Lipschitz
domains, for all finite growth exponents $r\ge2$ and all the prescribed mixed
pressure and flux data in their original Lebesgue and Sobolev spaces.
The solution can be chosen with speed at least one almost everywhere.
Finite divergence-free and gradient oscillations use the free unit sphere,
while monotone auxiliary equations confine the only constitutive defects
to bounded values. Comparing their exact solutions with the oscillatory
targets gives summable corrections in the conjugate norms, so the argument
also applies to rough data.
\end{abstract}

\section{Introduction}\label{sec:intro}
Adaptive porous-flow models select different drag laws according to the
local fluid speed. In the stationary mixed formulation, a drag field balances
the pressure gradient and an external force, while the velocity satisfies mass
conservation. Fumagalli and Patacchini introduced a model in which two radial
laws are selected below and above a speed threshold \cite{FP22}. Such a
switch permits the constitutive response to change across the medium rather
than fixing one law throughout the domain.

The sign of the drag jump changes the existence argument. The variational
theory in \cite[Theorems 4.10--4.11]{FP22} treats the nonnegative jump in
arbitrary dimension, whereas its negative-jump result is restricted to one
dimension. The related monotone framework in
\cite[Theorem 3.3, Remark 3.4 and Corollary 4.5]{FPpreprint} makes the same
distinction between general monotone laws and higher-dimensional nonmonotone
existence. A negative jump destroys convexity of the switched dissipation.
In several dimensions the mass constraint also permits nontrivial oscillations,
so minimization cannot be transferred from the one-dimensional argument.

There is an important distinction in the original model: the drag at the
threshold is completely free. Any vector $L^s$ field may be used on the
unit-speed set, not only a radial branch value or an energy subgradient.
As observed in \cite[Remark 3.6]{FP22}, weak solutions with this transition
law need not be minimizers of the underlying energy. We address precisely
that weak existence question, with the original mixed pressure and normal-flux
conditions and with no extra boundedness assumption on the data.

Our result is that every such datum in dimension $d\ge2$, with finite
growth exponent $r\ge2$ and conjugate exponent $s=r/(r-1)$, admits a
velocity in $L^r$, a pressure in the prescribed affine $W^{1,s}$ space,
and a drag in $L^s$ satisfying the constitutive and weak mass equations.
Moreover, the velocity can be chosen with speed at least one almost everywhere.
\Cref{thm:existence} states the complete hypotheses. Its inputs and solution
concept are those of the weak-existence question in \cite{AIM527}; the extra
speed bound is an output. Thus the theorem answers that question for every
original datum, without requiring minimization of the switched energy.

The free unit sphere permits finite splits of an arbitrary velocity--drag
pair into graph states, with divergence-free velocity differences and gradient
drag differences. Near the graph these splits have small mixed moments in
the two conjugate norms. Local finite-potential constructions of this kind
are part of convex integration \cite{ST21}. We give an explicit version
on measurable sets that preserves the global pressure and both original
boundary conditions.

The main difficulty is using such finite constructions with rough data:
their bounds depend on the pair values, which need not be uniformly bounded.
We first extend the high branch to a strongly monotone radial law, then add
a bounded source multiplied by a velocity cutoff. Every resulting exact
solution already lies on the high graph at large speeds; all its remaining
defects occur in a bounded pair range. Finite oscillations can therefore
reduce those defects. A uniform monotonicity comparison turns the approximate
affine fields into exact solutions of another cutoff equation. Iterating
with summable norm corrections gives a strong limit in the original spaces.

The proof uses the threshold's unrestricted drag essentially. It does not
give a selected transition law, uniqueness, or an energy-minimizing solution.
\Cref{sec:aux} gives the weak formulation and auxiliary equations;
\cref{sec:oscillations} proves the finite constructions;
\cref{sec:proof} completes the strong-limit argument. Equalities of fields
are understood almost everywhere and divergence identities distributionally.
Vector norms are Euclidean and $\norm{\cdot}{p}$ denotes the $L^p(\Omega)$ norm.

\section{Problem and auxiliary equations}\label{sec:aux}

\subsection{The weak problem}
Let $d\ge2$ and let $\Omega\subset\R^d$ be a bounded connected Lipschitz
domain. Its boundary is partitioned, up to surface measure zero, into relatively
open sets $\Sigma_p$ and $\Sigma_v$. Fix
\begin{equation}\label{eq:exponents}
 2\le r<\infty,\qquad s=\frac r{r-1}.
\end{equation}
The two coefficients $\phi_1:[0,1]\to[0,\infty)$ and
$\phi_2:[1,\infty)\to[0,\infty)$ are continuous and nondecreasing. Assume
that for some $c,C>0$,
\begin{equation}\label{eq:law-assumptions}
 c a^{(r-2)/2}\le\phi_2(a)\le C(1+a^{(r-2)/2})\quad(a\ge1),
 \qquad 0<\phi_2(1)<\phi_1(1).
\end{equation}
The negative jump is the last strict inequality. For
$u\in L^r(\Omega;\R^d)$ the permitted drag fields form the set
\begin{equation}\label{eq:original-law}
 \Lambda(u)=\left\{
 \phi_1(|u|^2)u\ind_{\{|u|<1\}}
 +\phi_2(|u|^2)u\ind_{\{|u|>1\}}
 +h\ind_{\{|u|=1\}}:
 h\in L^s(\Omega;\R^d)\right\}.
\end{equation}
In particular the unit-speed drag is an arbitrary vector field in $L^s$.
Each branch coefficient is evaluated only on its indicated speed region.

The data are
\[
 q\in L^r(\Omega),\qquad f\in L^s(\Omega;\R^d),\qquad
 u_0\in L^r(\Sigma_v).
\]
When $\Sigma_p$ has positive surface measure, prescribe an extension
$P\in W^{1,s}(\Omega)$ and set
\[
 V=\{\psi\in W^{1,s}(\Omega):\Tr\psi=0\text{ on }\Sigma_p\}.
\]
When $\Sigma_p=\varnothing$, use
\begin{equation}\label{eq:pure-flux}
 P=0,\qquad V=\left\{\psi\in W^{1,s}(\Omega):\int_\Omega\psi=0\right\},
 \qquad \int_\Omega q=\int_{\partial\Omega}u_0.
\end{equation}
These are the only cases: a nonempty relatively open part of a Lipschitz
boundary has positive surface measure. Boundary integrals below use the
Sobolev trace and surface measure.

A weak solution is a triple $u\in L^r(\Omega;\R^d)$, $p\in P+V$ and
$z\in L^s(\Omega;\R^d)$ such that
\begin{equation}\label{eq:weak-problem}
 z\in\Lambda(u),\qquad z=f-\nabla p,\qquad
 \pair{u}{\nabla\psi}=-\int_\Omega q\psi
      +\int_{\Sigma_v}u_0\Tr\psi\quad(\psi\in V).
\end{equation}
The velocity need not have a classical boundary trace; the last identity
specifies the normal-flux condition at its actual regularity.

\begin{theorem}\label{thm:existence}
For every domain, boundary partition, coefficients and datum satisfying
\eqref{eq:exponents}--\eqref{eq:pure-flux} and the conditions above,
\eqref{eq:weak-problem} has a weak solution. It can be chosen with
$|u|\ge1$ almost everywhere in $\Omega$.
\end{theorem}

The additional speed bound is a property of the constructed solution. It
removes the low branch from the proof; it imposes no restriction on the data.
The theorem concerns the free transition law \eqref{eq:original-law}. It
does not assert uniqueness or minimization of the switched dissipation.

\subsection{The affine constraints and a monotone extension}
Write $X=L^r(\Omega;\R^d)$ and define
\begin{align*}
 \ell(\psi)&=-\int_\Omega q\psi+\int_{\Sigma_v}u_0\Tr\psi,\\
 \mathcal C&=\{u\in X:\pair{u}{\nabla\psi}=\ell(\psi)\quad\forall\psi\in V\},\\
 K&=\{v\in X:\pair{v}{\nabla\psi}=0\quad\forall\psi\in V\},
 \qquad Y=\nabla V\subset L^s(\Omega;\R^d).
\end{align*}
The following linear facts will recover pressure after the velocity has
been constructed. They do not depend on either nonlinear branch.

\begin{proposition}[Linear pressure and mass constraints]\label{prop:linear}
The gradient norm on $V$ is equivalent to its $W^{1,s}$ norm.
The set $\mathcal C$ is a nonempty weakly closed affine translate of $K$,
and
\[
 K^\perp=Y=\nabla V.
\]
The gradient map $V\to Y$ has a bounded inverse, with $Y$ closed in $L^s$.
\end{proposition}
\begin{proof}
These are the linear statements of \cite[Lemmas 4.1--4.2, pp.\ 573--574]{FP22},
using the homogeneous pressure space defined on p.\ 572. The source's
pressure space is our $V$ and its velocity kernel is our $K$. Its operator
$B(u)(\psi)=\pair{u}{\nabla\psi}$ induces an isomorphism
$X/K\to V^*$, and its adjoint is an isomorphism $V\to K^\perp$.
The Poincare and trace bounds on this space make $\ell\in V^*$,
so $B(u)=\ell$ has a representative in $X$.
The adjoint is exactly the gradient under $L^r$--$L^s$ duality.
Its image is closed because annihilation of each member of $K$ passes
through an $L^s$ limit. Finally each defining gradient pairing for
$\mathcal C$ is weakly continuous. The source's linear proof uses the
trace constraints and boundary integrals only up to null sets, so the
present boundary-partition convention gives the same statements.
\end{proof}

Set $\beta=\phi_2(1)>0$ and extend the high law continuously below its
threshold:
\begin{equation}\label{eq:A}
 \alpha(t)=\begin{cases}\beta,&0\le t\le1,\\\phi_2(t^2),&t>1,\end{cases}
 \qquad A(a)=\alpha(|a|)a\quad(a\in\R^d).
\end{equation}
This extension permits exact monotone auxiliary equations for the original
data, without changing the weak affine constraints.

\clearpage
\begin{lemma}\label{lem:monotonicity}
There is $\gamma>0$, depending only on $\beta,c,r$, such that
for all $a,b\in\R^d$,
\begin{equation}\label{eq:monotonicity}
 (A(a)-A(b))\cdot(a-b)\ge\beta|a-b|^2,
 \qquad (A(a)-A(b))\cdot(a-b)\ge\gamma|a-b|^r.
\end{equation}
Also $A:L^r(\Omega;\R^d)\to L^s(\Omega;\R^d)$ is continuous.
More generally, if $a_n,c_n$ have uniformly integrable $r$th powers and
$\norm{a_n-c_n}{r}\to0$, then $\norm{A(a_n)-A(c_n)}{s}\to0$.
\end{lemma}
\begin{proof}
The coefficient $g=\alpha-\beta$ is nonnegative and nondecreasing.
If $|a|\ge|b|$, then
\[
 (g(|a|)a-g(|b|)b)\cdot(a-b)
 =g(|b|)|a-b|^2+(g(|a|)-g(|b|))(|a|^2-a\cdot b)\ge0.
\]
This proves the quadratic bound. Put $k_0=\min(\beta,c)$;
$\alpha(t)\ge k_0t^{r-2}$. Let $R_0=\max(|a|,|b|)$ and
$t_0=\min(|a|,|b|)$. If $t_0\ge R_0/2$, the same expansion gives
\[
 (A(a)-A(b))\cdot(a-b)
 \ge k_0(R_0/2)^{r-2}|a-b|^2
 \ge k_0 2^{4-2r}|a-b|^r.
\]
If $t_0<R_0/2$, use $a\cdot b\le R_0t_0$ to obtain
\begin{align*}
 (A(a)-A(b))\cdot(a-b)
 &\ge(R_0-t_0)(\alpha(R_0)R_0-\alpha(t_0)t_0)\\
 &\ge\alpha(R_0)(R_0-t_0)^2
 \ge(k_0/4)R_0^r
 \ge(k_0/4)(2/3)^r|a-b|^r.
\end{align*}
Thus one can take
$\gamma=k_0\min(2^{4-2r},\tfrac14(\tfrac23)^r)$.
The case $a=b=0$ is immediate; no division by a speed is required.

The growth assumptions imply
\begin{equation}\label{eq:growth}
 |A(a)|^s\le C_A(1+|a|^r).
\end{equation}
For the stated families, bounded $L^r$ norms make their large-value sets
small in measure. On a fixed velocity ball $A$ is uniformly continuous.
Together with $a_n-c_n\to0$ in measure this gives
$A(a_n)-A(c_n)\to0$ in measure. Their $s$th-power differences are
uniformly integrable by \eqref{eq:growth}, hence converge to zero in $L^s$.
Indeed, split their integral into a set where the difference exceeds a
fixed small number and its complement; the first set has vanishing measure
and the second contribution is bounded by that number to the power $s$
times $|\Omega|$.
Strong $L^r$ convergence makes the $r$th powers uniformly integrable,
which proves continuity at a fixed field. To see the stability under
small errors explicitly, write $c_n=a_n+e_n$ with $e_n\to0$ in $L^r$.
Then $|c_n|^r\le2^{r-1}(|a_n|^r+|e_n|^r)$. The total integrals
of $|e_n|^r$ are arbitrarily small after a finite prefix, and the prefix
has absolutely continuous integrals. This proves uniform integrability.
In particular, it holds for a fixed rough field plus uniformly bounded
increments and for all families close to it in $L^r$.
\end{proof}

\begin{lemma}\label{lem:auxiliary}
For every $F\in L^s(\Omega;\R^d)$ there is a unique pair
$u\in\mathcal C$, $p\in P+V$ satisfying $A(u)=F-\nabla p$.
\end{lemma}
\begin{proof}
The convex $C^1$ function
\[
 J(a)=\begin{cases}
 \beta|a|^2/2,&|a|\le1,\\
 \beta/2+\displaystyle\int_1^{|a|}\phi_2(t^2)t\,dt,&|a|>1
 \end{cases}
\]
has gradient $A$ and bounds
$(c/r)|a|^r-C_J\le J(a)\le C_J(1+|a|^r)$.
Minimize
\[
 \int_\Omega J(u)-\pair{F-\nabla P}{u}
 \quad\text{over }u\in\mathcal C.
\]
Coercivity and Holder's inequality bound a minimizing sequence in $L^r$.
Reflexivity and weak sequential compactness give a weak limit in
$\mathcal C$; see \cite[Theorems 3.18 and 4.10]{Brezis}.
The supporting-plane inequality at this limit,
$J(v)\ge J(u)+A(u)\cdot(v-u)$, proves weak lower semicontinuity after
integration, since $A(u)\in L^s$. The minimum is therefore attained.

For $w\in K$, differentiating along $u+tw$ gives
$\pair{A(u)-F+\nabla P}{w}=0$. The segment identity for $J$ and
\cref{lem:monotonicity} justify the derivative in the dual pairing;
its integrand is bounded by
$C(1+|u|^{r-1}+|w|^{r-1})|w|$, which is integrable, so the exchange
of the segment and spatial integrals is justified. By \cref{prop:linear}
there is $v\in V$ with $\nabla v=F-\nabla P-A(u)$.
Then $p=P+v$ has the required sign and trace convention.
For two solutions, their velocity difference is in $K$ and their
pressure difference in $V$. Testing the equation difference by the
velocity difference cancels pressure; \eqref{eq:monotonicity} and
\cref{prop:linear} give uniqueness.
\end{proof}

\subsection{Cutoff equations and bounded defects}
We will alter drag on bounded velocity sets while leaving large velocities
on the high graph. For $R\ge1$ use the $1/R$-Lipschitz function
\begin{equation}\label{eq:cutoff}
 \chi_R(a)=\begin{cases}
 1,&|a|\le R,\\2-|a|/R,&R<|a|<2R,\\0,&|a|\ge2R.
 \end{cases}
\end{equation}
For a bounded measurable source $b:\Omega\to\R^d$, set
$F_{b,R}(x,a)=A(a)+b(x)\chi_R(a)$.

\begin{lemma}\label{lem:cutoff}
If $M=\norm{b}{\infty}$ and $M/R\le\beta/4$, the original affine
constraints have an exact solution
\begin{equation}\label{eq:cutoff-solution}
 u\in\mathcal C,\quad p\in P+V,\quad
 z=F_{b,R}(x,u)=f-\nabla p.
\end{equation}
Moreover, if an affine pair $(\widetilde u,\widetilde z)$ satisfies
$\widetilde z=F_{b,R}(x,\widetilde u)+d$, $d\in L^s$, then
\begin{equation}\label{eq:comparison}
 \norm{u-\widetilde u}{r}^{r-1}\le\frac2\gamma\norm{d}{s}.
\end{equation}
Here an affine pair means $\widetilde u\in\mathcal C$ and
$\widetilde z=f-\nabla\widetilde p$ with $\widetilde p\in P+V$.
\end{lemma}
\begin{proof}
Average the two bounds in \eqref{eq:monotonicity} and subtract the
cutoff error $M|a-c|^2/R$:
\begin{equation}\label{eq:cutoff-monotonicity}
 (F_{b,R}(x,a)-F_{b,R}(x,c))\cdot(a-c)
 \ge\frac\gamma2|a-c|^r+\frac\beta4|a-c|^2.
\end{equation}
To obtain existence, for $v\in L^2$ let $T(v)$ be the velocity from
\cref{lem:auxiliary} with force $f-b\chi_R(v)$. This belongs to
$L^r\subset L^2$. Comparing two such equations, the pressure term
vanishes and the quadratic bound gives
\[
 \norm{T(v)-T(w)}2\le\frac{M}{\beta R}\norm{v-w}2
 \le\frac14\norm{v-w}2.
\]
Iteration is Cauchy in the complete space $L^2$; its limit is a fixed
point by this Lipschitz bound. As a value of $T$ it belongs to $L^r$,
and its auxiliary pressure gives \eqref{eq:cutoff-solution}.
For comparison, both velocity fields have the same affine mass data,
so their drag difference pairs to zero with their velocity difference.
Using \eqref{eq:cutoff-monotonicity} and Holder gives
$(\gamma/2)\norm{u-\widetilde u}{r}^r
 \le\norm{d}{s}\norm{u-\widetilde u}{r}$, proving
\eqref{eq:comparison}, including the zero-difference case.
\end{proof}

The closed part of the original graph which we shall use is
\begin{equation}\label{eq:graph}
 \Gamma=\{(a,z):|a|=1,\ z\in\R^d\}
 \cup\{(a,A(a)):|a|>1\}.
\end{equation}
Its continuous nonnegative defect is
\begin{equation}\label{eq:defect}
 G(a,z)=\begin{cases}
 1-|a|,&|a|\le1,\\
 \min(|a|-1,|z-A(a)|),&|a|>1.
 \end{cases}
\end{equation}
Continuity at the sphere follows from $G\le\big||a|-1\big|$ there,
and $G^{-1}(0)=\Gamma$.
For a solution of \eqref{eq:cutoff-solution}, the residual
$z-A(u)=b\chi_R(u)$ is bounded and vanishes at $|u|\ge2R$.
Consequently
\begin{equation}\label{eq:defect-bound}
 G(u,z)>0\ \Longrightarrow\ |u|<2R,\quad
 |z|\le\max_{|a|\le2R}|A(a)|+M.
\end{equation}
This is the boundedness used below. The solution itself may have
unbounded velocity and drag.

\section{Finite oscillatory constructions}\label{sec:oscillations}

\subsection{Wave connections and the full finite hull}
For pairs $U_i=(a_i,z_i)\in\R^d\times\R^d$ the relation
\begin{equation}\label{eq:wave}
 (a_1-a_2)\cdot(z_1-z_2)=0
\end{equation}
permits a divergence-free velocity oscillation and a gradient drag
oscillation with one common normal. Indeed there is a unit vector
$\xi$ such that $a_1-a_2$ is perpendicular to $\xi$ and
$z_1-z_2$ is parallel to $\xi$: use its direction if the drag
difference is nonzero, and otherwise choose a perpendicular direction
to the velocity difference. This uses $d\ge2$.

A finite laminate here means a finite probability measure obtained from
a single mean pair by finitely many binary splits
$U=\theta U_1+(1-\theta)U_2$, $0\le\theta\le1$, each satisfying
\eqref{eq:wave}. Zero-weight states can be omitted. We first establish
that every pair can be split into states of the graph \eqref{eq:graph}.

\begin{lemma}\label{lem:hull}
Every pair $(a,z)\in\R^{2d}$ is the mean of a finite laminate
supported in $\Gamma$. On each set $|a|,|z|\le D$ the atoms can
be chosen in a uniformly bounded set, with uniformly bounded split depth.
\end{lemma}
\begin{proof}
For $|a|=\rho\le1$, write $a=\rho e$ and choose a unit vector
$b\perp e$. At $a=0$ choose any orthonormal $e,b$.
The two unit states
\[
 \big(\rho e\pm\sqrt{1-\rho^2}\,b,z\big)
\]
with equal weights have mean $(a,z)$ and satisfy \eqref{eq:wave}.
When $\rho=1$ the original graph state suffices.

For $|a|>1$, choose a unit $b\perp a$, put $\delta=A(a)-z$,
and choose $T>|\delta|+|a\cdot\delta|+1$. For $\sigma\in\{-1,1\}$ set
\[
 z_\sigma=z+\sigma Tb,\qquad
 e_\sigma=\frac{A(a)-z_\sigma}{|A(a)-z_\sigma|},\qquad
 d_\sigma=a-e_\sigma.
\]
The denominator is nonzero by the choice of $T$. Along the ray
$y_\sigma(\lambda)=e_\sigma+\lambda d_\sigma$, $\lambda\ge1$,
the speed is at least $|a|>1$ and $d_\sigma\cdot y_\sigma(\lambda)>0$.
Consider
\[
 F_\sigma(\lambda)=d_\sigma\cdot(A(y_\sigma(\lambda))-z_\sigma).
\]
At the initial point,
\[
 F_\sigma(1)=a\cdot\delta-|\delta-\sigma Tb|<-1.
\]
The quadratic bound in \eqref{eq:monotonicity} gives, for $\lambda>\mu\ge1$,
\[
 F_\sigma(\lambda)-F_\sigma(\mu)
 \ge\beta(\lambda-\mu)|d_\sigma|^2.
\]
Thus the continuous function has a unique zero $\lambda_\sigma>1$.
Define
\[
 a_\sigma^h=e_\sigma+\lambda_\sigma d_\sigma,\qquad
 z_\sigma^u=
 \frac{\lambda_\sigma z_\sigma-A(a_\sigma^h)}{\lambda_\sigma-1}.
\]
The unit state $(e_\sigma,z_\sigma^u)$ and the high state
$(a_\sigma^h,A(a_\sigma^h))$, with weights
$(\lambda_\sigma-1)/\lambda_\sigma$ and $1/\lambda_\sigma$,
have mean $(a,z_\sigma)$. Their velocity difference is
$-\lambda_\sigma d_\sigma$, their drag difference is
$\lambda_\sigma(z_\sigma-A(a_\sigma^h))/(\lambda_\sigma-1)$,
and their scalar product is zero by the root equation.
Finally $(a,z_+)$ and $(a,z_-)$ have the same velocity and mean
$(a,z)$ with equal weights. This gives a two-level laminate.

For the uniform bounds, choose $T$ with a common strict margin on
the compact input range, including its limiting sphere $|a|=1$.
Let $D_\delta$ bound $|\delta|$. The transverse component gives
\[
 |d_\sigma|\ge|e_\sigma\cdot b|
 \ge\frac1{T+D_\delta}=:m>0.
\]
If $C_F$ bounds $|F_\sigma(1)|$, the preceding monotonicity gives
$\lambda_\sigma\le1+C_F/(\beta m^2)$.
On this bounded range the parameters and $\lambda$ form a compact
set. Uniform continuity of $F_\sigma$, together with
$F_\sigma(1)<-1$, keeps $\lambda_\sigma-1$ uniformly positive.
It therefore also bounds $z_\sigma^u$. The high atoms and the
subunit construction are bounded as well. No continuous choice of the
perpendicular vector is needed.
\end{proof}

\subsection{Mixed moment bounds near the graph}
Fullness of the finite hull alone does not give small corrections near
$\Gamma$. The following estimate supplies that property in the two
conjugate norms of the weak problem.

\begin{lemma}\label{lem:tight}
For every $D\ge1$ there are finite $B_D,C_D$ such that every
$U=(a,z)$ with $|a|,|z|\le D$ is the mean of a finite laminate
$\nu_U$ supported in $\Gamma\cap\{|a'|,|z'|\le B_D\}$, with
\begin{equation}\label{eq:moment}
 \int_{\R^{2d}}\big(|a'-a|^r+|z'-z|^s\big)\,d\nu_U(a',z')
 \le C_D G(a,z).
\end{equation}
The number of atoms and split depth can be bounded independently of $U$.
\end{lemma}
\begin{proof}
Inside the unit ball, the construction in \cref{lem:hull} has velocity
moment $(1-\rho^2)^{r/2}\le2^{r/2}(1-\rho)$ and zero drag moment.

For $a=(1+\varepsilon)e$, $0<\varepsilon<1/4$, choose a unit
$b\perp e$, a fixed $T>0$, and put
\[
 \theta=2\varepsilon,\qquad
 a_+=2e,\qquad
 a_-=\rho e,\quad \rho=\frac{1-3\varepsilon}{1-2\varepsilon},
 \qquad z_+=z+Tb,\quad
 z_-=z-\frac{\theta}{1-\theta}Tb.
\]
The weights $\theta,1-\theta$ give mean $(a,z)$, and the two
differences in this split are perpendicular. Replace the plus pair by
\cref{lem:hull}'s laminate and the minus pair by two unit states
with common drag $z_-$. The plus inputs stay in a fixed compact set,
so their atoms have a uniform bound and contribute at most
$C_D\varepsilon$. For each minus unit state,
\[
 |a'-a|^2\le9\varepsilon^2+4\varepsilon,\qquad
 |z_--z|\le4T\varepsilon.
\]
Here $1-\rho=\varepsilon/(1-2\varepsilon)\le2\varepsilon$.
Its mixed moment is at most
$C(\varepsilon^{r/2}+\varepsilon^s)\le C\varepsilon$, since
$r\ge2$ and $s>1$. This gives a bound by $|a|-1$ near the sphere.
When $|a|-1\ge1/4$, the bounded atoms from \cref{lem:hull} give
the same bound after enlargement of the constant.

For a second bound, suppose $|a|>1$ and
$\delta=|z-A(a)|\in(0,1/2)$. Split the pair into the exact high
state $(a,A(a))$ of weight $1-\delta$ and the mean pair
\[
 \left(a,A(a)+\frac{z-A(a)}\delta\right)
\]
of weight $\delta$. Their velocities agree, so the split is allowed.
The second input stays bounded; replace it using \cref{lem:hull}.
The first atom has drag deviation $\delta$, and the other bounded
atoms have total weight $\delta$. The mixed cost is at most
$C_D\delta$. For $\delta\ge1/2$ use the full bounded construction,
and for $\delta=0$ use the exact graph state. Selecting the cheaper
of this construction and the sphere construction proves \eqref{eq:moment}.
Each case uses a fixed finite binary tree, so the bounds on atom number
and depth are uniform.
\end{proof}

Enlarge the compact-set constants to nondecreasing majorants. Choose a
continuous nondecreasing $k:[1,\infty)\to[1,\infty)$ which dominates
them with one unit of radius slack; interpolation of a nondecreasing
sequence with $k(n)\ge C_{n+2}$ suffices. Set
\begin{equation}\label{eq:H}
 H(a,z)=k(1+|a|+|z|)G(a,z).
\end{equation}
Then $H\ge G$, its zero set is $\Gamma$, and at every pair one
of the laminates has mixed cost at most $H(a,z)$.
The function is bounded on compact pair sets. By \eqref{eq:defect-bound},
$H(u,z)$ is bounded for every cutoff solution, although no $r$-growth
bound for $H$ at arbitrary mismatched pairs is asserted.

\subsection{Measurable masks and exact affine perturbations}
The local finite-potential construction for laminates is familiar from
convex integration; see, for example, \cite[Proposition 4.5]{ST21}.
We give the required form directly, because the initial affine pair may
be rough and the set on which it is changed is merely measurable.

\begin{lemma}\label{lem:masks}
Let $U=(u,z)$ be an affine pair, and let $E$ be a measurable subset
of a compact set contained in $\Omega$, with $U$ bounded on $E$.
For every $\epsilon>0$ there are fields $W_n=(w_n,\omega_n)$ and
affine pairs $\widetilde W_n=(\widetilde w_n,\widetilde\omega_n)$ such that
\begin{align}
 &W_n=U\text{ off }E,\quad W_n-U\text{ uniformly bounded},\label{eq:target-support}\\
 &\int_E H(W_n)<\epsilon,\label{eq:target-H}\\
 &\limsup_n\big(\norm{w_n-u}{r}^r+\norm{\omega_n-z}{s}^s\big)
      \le\int_E H(U)+\epsilon,\label{eq:target-cost}\\
 &\norm{\widetilde w_n-w_n}{r}
   +\norm{\widetilde\omega_n-\omega_n}{s}\longrightarrow0.
   \label{eq:affine-approx}
\end{align}
The perturbations $\widetilde W_n-U$ are compactly supported
divergence-free/gradient pairs. For this fixed construction they also
have a uniform $L^\infty$ bound.
\end{lemma}
\begin{proof}
Partition $E$ into finitely many measurable value bins $E_j$, with
constants $C_j$ such that $|U-C_j|\le\delta$ on $E_j$.
Here and below a pair norm is $|(a,z)|=|a|+|z|$.
Choose $\delta\le1$, so all centers remain in one fixed compact
range. Choose the laminates of \cref{lem:tight} at those centers.
Their atoms have a common bound independent of the bin precision.
On a bin the ideal field takes values
\[
 U(x)+(a_{j,\ell},z_{j,\ell})-C_j,
\]
where $(a_{j,\ell},z_{j,\ell})$ is a graph atom. Uniform continuity
of $H$ on that compact range makes its integral on $E$ smaller
than $\epsilon$ by choosing $\delta$ small. It also makes
$\sum_j|E_j|H(C_j)\le\int_E H(U)+\epsilon$ with any prescribed
smaller error in place of $\epsilon$.

We realize the state fractions while approximating these ideal fields
by exact affine ones. For a binary split with difference
$q=(q_u,q_z)$ and weights $\theta,1-\theta$, choose a unit
$\xi$ with $q_u\cdot\xi=0$ and $q_z=t\xi$.
Let $h$ be a periodic profile taking values $1-\theta$ and
$-\theta$ on fractions $\theta$ and $1-\theta$ of each period.
Its mean is zero, so it has a bounded periodic Lipschitz primitive $L$.
For a measurable parent phase $F$, the ideal correction is
$\ind_F h(N\xi\cdot x)q$.

Approximate $\ind_F$ by smooth amplitudes $\eta$, between zero
and one, supported in a fixed compact interior neighborhood. Convolution
gives convergence in both $L^r$ and $L^s$; see
\cite[Theorem 4.22]{Brezis}. Set
$S=q_u\otimes\xi-\xi\otimes q_u$, so $S$ is skew and $S\xi=q_u$.
The potentials
\begin{equation}\label{eq:potentials}
 \Psi_N(x)=\frac1N S\eta(x)L(N\xi\cdot x),\qquad
 v_N(x)=\frac tN\eta(x)L(N\xi\cdot x)
\end{equation}
give
\[
 (\Div\Psi_N,\nabla v_N)
 =\eta h(N\xi\cdot x)q
 +\frac{L(N\xi\cdot x)}N(S\nabla\eta,t\nabla\eta).
\]
Row divergence is used for the matrix. First choose the smooth amplitude
with small mask error, and then choose $N$ so the last term is small
in $L^\infty$. Differentiation is also valid distributionally: $L$ is
continuous and piecewise affine, so integration by parts on its affine
intervals has cancelling interface terms. Fubini along the $\xi$
coordinate gives the displayed derivative identity.

The velocity perturbation is divergence-free because the potential is
skew. It pairs to zero with the gradient of every $\psi\in V$:
first use smooth tests, then the local Sobolev approximation of
\cite[Theorem 9.2]{Brezis} on the compact support, and pass by Holder.
The drag perturbation is a global gradient. Its pressure correction is
$-v_N$, with the mean restored by a constant in the pure-flux case.
Thus both original boundary conditions are preserved, also on domains
with holes.

Periodic averaging gives the required fractions on every measurable
parent phase. For a smooth compactly supported weight, integrate a
zero-mean profile's bounded primitive in direction $\xi$, obtaining
an error $O(N^{-1})$. Density in $L^1$ extends this to $\ind_F$;
see \cite[Theorem 4.12]{Brezis} and mollification. Apply this also
to each phase indicator minus its mean fraction.

Carry out the finite split trees in order. At index $n$, choose the
outer frequencies with errors below $1/n$. Their ideal phase masks
are then fixed measurable sets. For each inner split, smooth its
parent mask with error below $1/n$ in both norms and choose its
frequency so its fraction error and potential remainder are below
$1/n$. There are only finitely many splits. Their free corrections
sum to a uniformly bounded compactly supported perturbation, and
the accumulated mask and remainder errors tend to zero.
On the disjoint bins the corresponding ideal targets take the leaf
values above, so they have a uniform bound on $E$.

The conditional phase measures converge to the prescribed weights.
Consequently the ideal mixed increment cost converges to the sum of
the bin measures times the laminate moments, bounded by
$\sum_j|E_j|H(C_j)$. Choosing the bin error smaller than
$\epsilon$ proves \eqref{eq:target-cost}. The other properties
follow from the construction. All bounds here are for a fixed finite
collection of bins and trees; no measurable selection of a continuum
of laminates is needed.
\end{proof}

\section{Proof of existence}\label{sec:proof}

We now use the finite oscillations to improve exact solutions, keeping the
original data fixed. Call a pair a compact-source cutoff solution if it
satisfies \eqref{eq:cutoff-solution}, $b$ vanishes outside a compact subset
of $\Omega$, and $R\ge1$, $\norm{b}{\infty}/R\le\beta/4$.
The zero source is allowed. Recall that $H$ is the continuous weighted
defect \eqref{eq:H}; it is integrable on every such solution by
\eqref{eq:defect-bound}.

\begin{lemma}[Exact improvement]\label{lem:improvement}
For every compact-source cutoff solution $U=(u,z)$ and every $\eta>0$
there is another one, $U'=(u',z')$, with the identical original data and
\begin{equation}\label{eq:improvement}
 \int_\Omega H(U')<\eta,\qquad
 \norm{u'-u}{r}^r+\norm{z'-z}{s}^s\le2\int_\Omega H(U)+\eta.
\end{equation}
\end{lemma}
\begin{proof}
Since $H(U)$ is bounded, choose a compact interior set $K_0$ with
$\int_{\Omega\setminus K_0}H(U)<\eta/10$ and set
$E=K_0\cap\{G(U)>0\}$. Compact interior exhaustion and dominated
convergence justify the choice. The same bound holds outside $E$,
and $U$ is bounded on $E$ by \eqref{eq:defect-bound}.
Use \cref{lem:masks} with an error $\epsilon<\eta/10$, small enough
that the target $H$ integral on $E$ is also below $\eta/10$.
Write its ideal targets $W_n=(w_n,\omega_n)$ and affine approximants
$\widetilde W_n=(\widetilde w_n,\widetilde\omega_n)$.

The target residual
\[
 b_n(x)=\omega_n(x)-A(w_n(x))
\]
is bounded on $E$ by the common target range. Outside $E$ it is
the old residual $b\chi_R(u)$, hence is bounded and vanishes off the
old compact source support. Thus all $b_n$ vanish outside one compact
set, and $\norm{b_n}{\infty}\le M_*$ uniformly in $n$.
Choose $R_*$, independently of $n$, larger than $2R$ and the
target velocity bound on $E$, with $M_*/R_*\le\beta/4$.
At every point where $b_n\ne0$, $|w_n|<R_*$, so
\begin{equation}\label{eq:target-identity}
 \omega_n=A(w_n)+b_n\chi_{R_*}(w_n).
\end{equation}
Solve the exact cutoff problem of \cref{lem:cutoff} with this source,
obtaining $U_n^*=(u_n^*,z_n^*)$ and the original mass and pressure data.

For comparison with the affine approximant define
\[
 d_n=\widetilde\omega_n-A(\widetilde w_n)
          -b_n\chi_{R_*}(\widetilde w_n).
\]
Subtract \eqref{eq:target-identity}. The direct drag difference
tends to zero in $L^s$. The targets are the fixed rough $u$ plus
uniformly bounded increments, so their $r$th powers are uniformly
integrable; the $L^r$-close affine velocities have the same property.
\Cref{lem:monotonicity} therefore makes their $A$ difference tend to
zero in $L^s$. The cutoff difference tends to zero in $L^s$ by
its Lipschitz bound, the common $M_*$ and $L^r\subset L^s$ on
the bounded domain. Consequently $\norm{d_n}{s}\to0$.

The exact and approximate velocities lie in $\mathcal C$ and both
drag fields equal $f$ minus a gradient with pressure in $P+V$.
\Cref{lem:cutoff} yields
\[
 \norm{u_n^*-\widetilde w_n}{r}^{r-1}
 \le\frac2\gamma\norm{d_n}{s}\longrightarrow0.
\]
Uniform integrability passes to these exact velocities. The identity
\[
 z_n^*-\widetilde\omega_n
 =A(u_n^*)-A(\widetilde w_n)
  +b_n\big(\chi_{R_*}(u_n^*)-\chi_{R_*}(\widetilde w_n)\big)-d_n
\]
and \cref{lem:monotonicity} then give drag convergence in $L^s$.
Thus $U_n^*-W_n\to0$ in $L^r\times L^s$.

The $H$ integrals also become close, despite the possible rough tails.
All exact solutions have the common cutoff radius $R_*$ and source
bound $M_*$, so \eqref{eq:defect-bound} bounds $H(U_n^*)$ uniformly.
The targets have a uniform $H$ bound as well: off $E$ they are $U$,
and on $E$ they lie in a common compact range. They are tight in
measure, being the fixed pair $U$ plus uniformly bounded increments.
Discard a small-measure large-value set of $U$. On its complement
$W_n$ lies in a common compact range; outside another set of vanishing
measure, product-norm closeness places $U_n^*$ in a slightly larger
range. Uniform continuity of $H$ there proves
$H(U_n^*)-H(W_n)\to0$ in measure. Its uniform bound makes the
convergence $L^1$: split at a fixed small magnitude as in the proof
of \cref{lem:monotonicity}.
Only these exact fields and ideal targets enter this argument; no
integrability of $H(\widetilde W_n)$ is needed.

It follows that $\int_\Omega H(U_n^*)<\eta$ for all sufficiently
large $n$. Also product-norm closeness and the bounded target increment
norms imply
\[
 \limsup_n\big(\norm{u_n^*-u}{r}^r+\norm{z_n^*-z}{s}^s\big)
 \le\int_\Omega H(U)+\epsilon.
\]
A sufficiently large finite $n$ satisfies both inequalities in
\eqref{eq:improvement}, including when the original integral is zero.
Set $U'=U_n^*$. The new source and cutoff bounds may depend on this
step, but were fixed throughout its frequency limit.
\end{proof}

\begin{proof}[Proof of \cref{thm:existence}]
Start with the auxiliary solution $U^{(0)}=(u^{(0)},z^{(0)})$ from
\cref{lem:auxiliary} with $F=f$. It is a compact-source cutoff
solution with $b=0$. On $|u^{(0)}|>1$ its defect vanishes, and on
$|u^{(0)}|\le1$ its drag is $\beta u^{(0)}$, so
$\int_\Omega H(U^{(0)})<\infty$.
Apply \cref{lem:improvement} to $U^{(j)}$ with
$\eta_j=2^{-r(j+1)}$. Then, for $j\ge1$,
\begin{align*}
 \int_\Omega H(U^{(j)})&\le2^{-rj},\\
 \norm{u^{(j+1)}-u^{(j)}}{r}^r
  +\norm{z^{(j+1)}-z^{(j)}}{s}^s&\le3\,2^{-rj}.
\end{align*}
The first step has finite cost. The remaining velocity norm increments
are bounded by $3^{1/r}2^{-j}$, and the drag norm increments by
$3^{1/s}2^{-(r-1)j}$. Both series converge. Completeness gives
\[
 u^{(j)}\longrightarrow u\text{ in }L^r,\qquad
 z^{(j)}\longrightarrow z\text{ in }L^s.
\]
Each $z^{(j)}=f-\nabla p^{(j)}$ with $p^{(j)}\in P+V$.
The bounded inverse gradient of \cref{prop:linear} makes
$p^{(j)}-P$ Cauchy in $W^{1,s}$, giving $p\in P+V$ and
$z=f-\nabla p$. For every fixed $\psi\in V$, Holder's inequality
passes the mass identity to the limit. This preserves exactly the
original normal-flux functional and, in pure flux, the pressure mean.

An almost-everywhere convergent subsequence exists in the finitely
many pair components by \cite[Theorem 4.9]{Brezis}. Continuity of
$H$ and Fatou's lemma give
\[
 0\le\int_\Omega H(u,z)
 \le\liminf_j\int_\Omega H(U^{(j)})=0.
\]
Therefore $(u,z)\in\Gamma$ almost everywhere. In particular $|u|\ge1$,
and $z=\phi_2(|u|^2)u$ where $|u|>1$. Define the single global
$h=z\ind_{\{|u|=1\}}\in L^s$. The original law
\eqref{eq:original-law} then gives $z\in\Lambda(u)$, and all
conditions in \eqref{eq:weak-problem} hold.
\end{proof}

\clearpage
\begin{thebibliography}{9}
\bibitem{Brezis}
H. Brezis, \emph{Functional Analysis, Sobolev Spaces and Partial Differential
Equations}, Springer, 2011.
\href{https://doi.org/10.1007/978-0-387-70914-7}{doi:10.1007/978-0-387-70914-7}.

\bibitem{FP22}
A. Fumagalli and F. S. Patacchini,
Model adaptation for non-linear elliptic equations in mixed form:
existence of solutions and numerical strategies,
\emph{ESAIM Math. Model. Numer. Anal.} \textbf{56} (2022), 565--592.
\href{https://doi.org/10.1051/m2an/2022016}{doi:10.1051/m2an/2022016}.

\bibitem{FPpreprint}
A. Fumagalli and F. S. Patacchini,
Well-posedness and variational numerical scheme for an adaptive model in
highly heterogeneous porous media,
\href{https://arxiv.org/abs/2206.07970v1}{arXiv:2206.07970v1}, 2022.

\bibitem{ST21}
M. Sorella and R. Tione,
The four-state problem and convex integration for linear differential operators,
\href{https://arxiv.org/abs/2107.10785v1}{arXiv:2107.10785v1}, 2021.

\bibitem{AIM527}
MColbrook/AIM, \emph{527. Existence for adaptive porous flow with a negative
drag jump}, problem repository, revision
\texttt{c9929805cece044705abaaebe1cebddac1ba11a6}, 2026.
\href{https://github.com/MColbrook/AIM/blob/c9929805cece044705abaaebe1cebddac1ba11a6/problems/527-negative-jump-adaptive-flow.md}{Immutable problem statement}.
\end{thebibliography}

\section*{AI-use statement}
This proof and manuscript were developed with substantial assistance from
OpenAI Codex. AI assistance included mathematical exploration, drafting,
checking primary-source statements, self-review and manuscript preparation.
A separate Codex context reviewed a prepublication version of the manuscript
and its cited public sources. Both drafting and AI review used OpenAI systems.
No independent human review or formal proof verification is claimed.
The argument is submitted for independent mathematical review.
\end{document}
